% Contextual sources now cited in the supporting body argument.
% Bibliographic records have not all been independently full-text audited.
\paragraph{Evaluation environments, native labels and reliability.}
Meta-evaluation and benchmarking surveys emphasize the dependence of an agent score on task definition, judge, sampling protocol and failure semantics~\cite{ext003,ext006,ext029}.
Benchmark auditing and unified evaluation efforts make explicit the difference between task completion, robustness, blame attribution and failure reporting~\cite{ext030,ext031,ext036}.
Therefore this review's capability matrix recognizes a native decisive-step label only when the benchmark's own annotation protocol provides it~\cite{ext038,ext039,ext139}.

General-agent and collaboration benchmarks measure performance across diverse cooperation and tool-use settings~\cite{ext072,ext073,ext133}.
Web, enterprise, desktop and software-repository environments provide more specific tasks and stateful interaction requirements~\cite{ext134,ext135,ext136}.
These tasks can generate useful traces, but end-state task success cannot substitute for independently labeled agent-level or step-level causes~\cite{ext137,ext138}.

Post-failure reporting and agent-as-a-judge evaluation introduce complementary ways to assess explanation quality and agent behavior~\cite{ext063}.
However, plausible narrative explanations and automatic judge ratings are not automatically equivalent to a verified causal failure step~\cite{ext101}.
Separate reporting and attribution metrics are therefore needed when selecting benchmarks for diagnostic experiments.

